\begin{proof}
The sum map $m|_{X\times X}:X\times X\to P$ is surjective because two
translates of the ample divisor $X$ always meet.  Moreover,
\[
 m\bigl((\{0\}\times X)\cup(X\times\{0\})\bigr)=X.
\]
Hence, for $t\notin X$, the fibre lies entirely in
$X_{\reg}\times X_{\reg}$.  Generic smoothness for the restricted sum map
therefore gives smoothness of the general fibre.  Since $t\notin X$ and
$\Stab_P(X)=\{0\}$, the translate $t-X$ does not contain $X$; its restriction
to the integral Cohen--Macaulay surface $X$ is therefore an effective ample
Cartier divisor.  Hartshorne's connectedness theorem, applied on $X$, shows
that this divisor, namely $X\cap(t-X)$, is connected
\cite[Corollary~III.7.9]{Hartshorne1977}.
The involution $p\mapsto t-p$ preserves its complex orientation and hence
acts by $+1$ on its one-dimensional top cohomology.  Therefore
\begin{equation}
 \cH^{-2}(K*K)_-=0
 \quad\text{on a dense open subset of }P.
\label{eq:generic-alt-vanishing}
\end{equation}
\end{proof}

\begin{lemma}[Extraction of the divisor-supported summand]
\label[lemma]{lem:extract-ic}
In the decomposition-theorem splitting of $(K*K)_-$, the constant rank-one
local system of \cref{lem:split-fibre-top} is the restriction of an unshifted
$\IC_X$ summand.  It cannot arise from a shifted simple summand with support
$P$, or from a summand with any other support.
\end{lemma}

\begin{proof}
The complex $(K*K)_-$ is a direct summand of the projective direct image of
a pure polarizable Hodge module.  By the decomposition theorem and
semisimplicity \cite{BBD1982,Saito1988}, it is a direct sum of shifted simple
intersection complexes.  A summand with proper support containing the
generic point of $U^+$ must have support $X$.  Such a simple summand has the
form $\IC_X(\mathcal V)[r]$, where $\mathcal V$ is an irreducible local
system on a dense smooth open $X_{\mathcal V}\subset X_{\reg}$.  Its
restriction to $U^+\cap X_{\mathcal V}$ is $\mathcal V[2+r]$, so contribution
to ordinary degree $-2$ forces $r=0$.

A full-support summand cannot account for the line in
\eqref{eq:principal-alt-rank}.  Choose a general point $a\in U^+$ that lies
on no singular stratum of any full-support simple summand except the divisor
stratum $X_{\reg}$.  A small analytic neighbourhood of $a$ is stratified
homeomorphic to a product of a ball in $X_{\reg}$ with a normal disk
$\Delta$ to $X$.  Hence a full-support simple perverse summand is locally
\[
 \mathcal L[2]\boxtimes j_{\Delta,!*}\mathcal W[1],
 \qquad j_\Delta:\Delta^*\hookrightarrow\Delta,
\]
for local systems $\mathcal L$ and $\mathcal W$.

Write $T$ for the local monodromy of $\mathcal W$.  On a disk the middle
extension has the explicit stalk
\[
 \cH^k\!\left(i_0^*j_{\Delta,!*}\mathcal W[1]\right)
 =\begin{cases}
   \ker(T-1),&k=-1,\\
   0,&k\ne-1,
  \end{cases}
 \qquad i_0:\{0\}\hookrightarrow\Delta.
\]
Indeed $Rj_{\Delta,*}\mathcal W[1]$ has stalk cohomology
$\ker(T-1)$ in degree $-1$ and $\operatorname{coker}(T-1)$ in degree $0$;
the middle extension removes exactly the quotient supported at the origin.
It follows that an unshifted full-support perverse summand contributes along
$X$ only in ordinary degree $-3$.  To contribute to ordinary degree $-2$
it must occur with shift $[-1]$.  But on the dense smooth open of $P$ the
same shifted summand is a nonzero local system $\mathcal E[3][-1]$, hence
also contributes to ordinary degree $-2$ there, contradicting
\eqref{eq:generic-alt-vanishing}.  Thus no full-support summand contributes
to the rank-one line on $U^+$.

The two relevant shifts are summarized by
\[
\begin{array}{c|c|c}
\text{full-support summand}&\text{general point of }P&\text{general point of }X\\ \hline
\IC_P(\mathcal E)&-3&-3\\
\IC_P(\mathcal E)[-1]&-2&-2.
\end{array}
\]
Supports of dimension at most one do not meet a generic point of $U^+$, and
a different irreducible divisor cannot contain the nonempty open $U^+\subset
X$.

The local system $\mathcal V$ has rank one because its contribution to
\eqref{eq:principal-alt-rank} has rank one.  Replace $U^+$ by
$U^+\cap X_{\mathcal V}$.  Its restriction there is the constant local
system found in \cref{lem:split-fibre-top}.  The complement
$X_{\mathcal V}\setminus U^+$ is a proper complex algebraic subset and
therefore has real codimension at least two.  Every loop in the smooth
variety $X_{\mathcal V}$ can be perturbed off this complement, so the
inclusion induces a surjection
\[
 \pi_1(U^+)\twoheadrightarrow\pi_1(X_{\mathcal V}).
\]
The monodromy representation of $\mathcal V$ is trivial on the source and
hence on $\pi_1(X_{\mathcal V})$.  Thus $\mathcal V$ is trivial and the
summand is $\IC_X$ itself.
\end{proof}

\begin{proposition}[Alternating support detector]
\label[proposition]{prop:alt-support}
In $\operatorname{Perv}(P,\QQ)$, $K=\IC_X$ is a direct summand of
${}^pH^0((K*K)_-)$.  After tensoring with $\CC$, its image in
$\overline{\operatorname{Perv}}(P)$ is nonzero and is a direct summand of
$\Alt^2(K)$.
\end{proposition}

\begin{proof}
The first assertion is \cref{lem:extract-ic}.  Apply the quotient functor to
$K\otimes_{\QQ}\CC$.  It is perverse $t$-exact, convolution is $t$-exact in
$\overline D(P)$, and all nonzero perverse-cohomological degrees of $K*K$ map
to zero.  Hence the
image of ${}^pH^0((K*K)_-)$ is the anti-invariant summand of the tensor square
of the image of $K$.  Because the support dimension is even, this is
$\Alt^2(K)$ under the fibre functor.  The exact quotient functor preserves
the direct summand supplied by \cref{lem:extract-ic}; that summand is nonzero
because $\chi(K)=7$.
\end{proof}

\section{Identification of the full Tannaka group}

Let
\[
 G=G(K),\qquad V=\omega(K).
\]
By \eqref{eq:chi-seven}, $\dim V=7$.  The self-duality
\eqref{eq:self-duality} gives a perfect evaluation pairing
$K*K\to\delta_0$.  The chosen symmetry
is centred at the origin, so no translated skyscraper occurs, and
$\dim X=2$ makes the pairing symmetric.  Therefore the full group embeds
faithfully as
\begin{equation}
 G\hookrightarrow \mathrm{O}(V).
\label{eq:orthogonal}
\end{equation}

The simple perverse sheaf $K$ is nondivisible: a translation fixing it must
fix its support, and \eqref{eq:stabilizer} then makes the translation zero.
By \cite[Proposition~3.3]{JKLM2025},
\begin{equation}
 V|_{G^0}\quad\text{is irreducible}.
\label{eq:connected-irreducible}
\end{equation}
This result does not require the ambient abelian variety to be simple.

Choose $\alpha\in\Omega$ as in \cref{prop:hodge-vector}.  By
\cref{lem:tensor-compatible-twist}, $H^0(-\otimes\mathcal L_\alpha)$ is a
tensor fibre functor on the category generated by $K$.  Since any two fibre
functors over $\CC$ are noncanonically isomorphic, it may be identified
with $\omega$.  Let $G_P=G$, and let
$G_H$ be the Tannaka group of the pure Hodge module $\IC_X^H$, acting
faithfully on the same complex vector space $V$.
By \cite[Corollary~3.2]{KLM2026}, the Hodge decomposition is defined by a
cocharacter $h:\GGm^2\to G_H$.  The comparison theorem
\cite[Corollary~2.5]{KLM2026} gives
\[
 G_P\mathrel{\triangleleft}G_H,
 \qquad
 [G_P^0,G_P^0]=[G_H^0,G_H^0]=:G^*.
\]
Since $G_P^0\subset G_H^0$, equation
\eqref{eq:connected-irreducible} implies that $V|_{G_H^0}$ is irreducible.
Thus the connected centre $Z_H^0$ acts by scalar matrices.  As $G_H^0$ is
reductive, multiplication $Z_H^0\times G^*\to G_H^0$ is surjective with
finite kernel.

